%=======================================================================
%  Paper V — Section 7 (Failure modes)   FINAL
%
%  Requires the macros listed at the head of Section 2.
%
%  PRINCIPAL CORRECTION IN THIS REVISION -----------------------------
%  Section 7.3 fixes the region fraction at the misplaced Eros plane at
%  f = 0.8844 (complement 11.56 %), recovered exactly from the reported
%  rho_rest = 19,259 at rho_R = 500. This overturns the claim made in
%  Section 4.4 of an earlier draft that the displacement statistic S
%  reaches 19 at that plane: the correct value is 2.86, BELOW the 5.8
%  returned by Itokawa, so S ranks the artefact correctly. What fails
%  instead is the amplification factor, |A| = 9.12 < 12.0, which ranks
%  the artefact as the stronger signal of the two. Section 4.4 and
%  Table 4.2 must be revised accordingly.
%
%  Superseded values:
%      "2.1e-14 in relative terms" volume closure -> 2.0e-14 per cent
%      "S = 19 at the misplaced plane" -> 2.86 (Section 4.4)
%      complement-too-small "survives: background density"
%                                      -> destroys it; 3.83 % vs 1.38 %
%      "7.5 kg/m3 per metre"           -> 10.0 at the reference plane,
%                                         and 0.28 for rho_rest
%      3700/4125 as a flatness demonstration
%                                      -> a sub-grid location of the
%                                         minimiser at rho_R ~ 3912
%
%  All values confirmed against paperV_final.json and paper4_recheck.json.
%  Resolved in this pass: closure residual is a percentage, not a relative
%  value; the "-355/+1114" pair conflated the 0.60 % and 5.9 % levels;
%  the Eros envelope is 6.03 % (quadric) or 5.40 % (clustering); the window
%  diagnostic reproduces 10.3 % once the bounds are applied to both
%  densities, and Itokawa recomputes to 87.6 %; the Arrokoth crossover with
%  the buried facets excluded lies within 1 kg/m3 of 250.
%=======================================================================

\section{Failure modes}
\label{sec:failures}

Most of what went wrong in this work went wrong silently. A region can
enclose the wrong solid and still report an exact volume; a minimum can be
pinned to the edge of a scan and still report an improvement; a set of
facets can lie inside the body and still enter a surface average. This
section catalogues the seven failures encountered, with the quantity each
destroys and the check that catches it, because a reader applying the
method needs the failure list more than the result list.

Two of them are recorded not because they are subtle but because they were
ours: the wrong solid of Section~\ref{sec:methods:cone} survived a
volume-closure test, and the buried facets of
Section~\ref{sec:methods:shapes} survived every geometric check the
implementation then had. The others are properties of the method that any
implementation will meet.

\begin{table}[htbp]
\centering
\scriptsize
\setlength{\tabcolsep}{3.5pt}
\caption{The seven failure modes. ``Destroys'' names the quantity that
becomes untrustworthy, and is stated in terms of the quantities
Section~\ref{sec:constrains} identifies as constrained ($\densrest$ and
$\Dmass$) rather than in terms of $\densR$, which the method does not
determine in any case. Note the asymmetry between the fourth and fifth
rows: a region that is too small degrades only $\densR$, while a
complement that is too small degrades $\densrest$, and it is the second
that matters.}
\label{tab:failures}
\setlength{\tabcolsep}{5pt}
\small
\begin{tabular}{p{2.3cm}p{2.6cm}p{3.6cm}p{2.4cm}p{3.0cm}}
\toprule
Failure & Encountered & Symptom & Destroys & Caught by \\
\midrule
Wrong solid built
 & Itokawa, all planes
 & volume exact, mesh watertight, contrast inverted
 & everything
 & vertex predicate; mass-balance audit \\
\addlinespace
Degenerate facets
 & Itokawa, $\xsplit = \SI{0.130}{\kilo\metre}$
 & NaN at every field point, not merely near the plane
 & everything
 & self-announcing \\
\addlinespace
Facets not on the surface
 & Arrokoth, two unjoined shells
 & $0.943\%$ of the area buried; a sign reversal appears that is not there
 & sign of $\Dmass$
 & per-component Euler check; containment test \\
\addlinespace
Region too small
 & Itokawa, $\SI{0.240}{\kilo\metre}$, $f = 5.0\%$
 & curvature interval widens from $\pm4.6\%$ to $\pm7.1\%$
 & $\densR$ only
 & curvature interval \\
\addlinespace
Complement too small
 & Eros, $\SI{-10.76}{\kilo\metre}$, complement $11.56\%$
 & curvature interval in $\densR$ \emph{narrows} to $\pm15$, which looks like precision
 & $\densrest$, at $3.83\%$ against $1.38\%$
 & interval propagated to $\densrest$, not read in $\densR$ \\
\addlinespace
Minimum at a scan edge
 & Arrokoth neck, $f = 1.0\%$
 & $\improve = 13.6\%$ from a density pinned to the scan floor
 & everything
 & boundary-minimum guard; specific improvement \\
\addlinespace
Plane misplaced
 & Eros, $\SI{-10.76}{\kilo\metre}$
 & $C = 0.876 \pm 0.039$, excluding unity at $3.2$ intervals; $\dCOM$ reversed and four times larger
 & everything
 & envelope comparison; plane provenance; admissible-scan fraction \\
\bottomrule
\end{tabular}
\end{table}

\subsection{Building the wrong solid}
\label{sec:failures:cone}

The first failure is the one described in Section~\ref{sec:methods:cone},
and it is placed first because it is the only one that produced a
publishable-looking wrong answer across every plane simultaneously.

Closing a clipped region by fanning its boundary edges to a single interior
point produces a watertight mesh whose divergence-theorem volume agrees
with the sum of its own tetrahedra to machine precision. The
volume-closure test then in use compared exactly those two quantities, and
reported a residual of $2.0\times10^{-14}$ per cent, that is
$2.0\times10^{-16}$ in relative terms---round-off, and nothing else. That
figure lies below the double-precision epsilon because the two half-space
meshes share most of their facets with the whole-body mesh, so the round-off
is common to both sides of the comparison and cancels
(Section~\ref{sec:methods:verify}). The solid enclosed was nevertheless a radial cone over
the surface patch rather than the plane-cut lobe, too large by a factor of
$1.600$ on a unit sphere and with its centroid displaced by $0.1125\,R$.
% Resolved from the code: density_estimate.py records
% complement_closure_pct = 100 x |V_R + V_complement - V_tot| / V_tot, so the
% archived 1.95e-14 is a percentage. Section 2.6 and this subsection now agree.

On Itokawa the consequence was a sign reversal in the recovered contrast,
because raising the density of a cone adds mass along a wedge reaching
back to the centre of the body rather than in the head. Two checks would
have caught it and neither was present: the vertex predicate of
Section~\ref{sec:methods:verify}, which the cone fails outright with a
violation of $0.5$ in units where the plane sits at $x = 0.5$; and the
mass-balance audit, which would have found
$\Dmass = \Ddens\,\VR$ and $\Dmass = (\densbulk-\densrest)\Vtot$
disagreeing by the same factor $1.600$, since the volume passed to the
mass constraint and the volume generating the potential were different
solids.

The general lesson is that a closure test comparing a mesh against its own
tetrahedra tests self-consistency, not correctness. Only a test that
appeals to the region's \emph{definition}---every vertex satisfying its
defining inequality---distinguishes the two constructions.

\subsection{Degenerate facets and total NaN}
\label{sec:failures:nan}

At $\xsplit = \SI{0.130}{\kilo\metre}$ on Itokawa a vertex of the source
mesh lies on the cutting plane. Without the classification rule of
Section~\ref{sec:methods:tol}, both edges incident on that vertex register
as crossings, two coincident intersection points are created, and
zero-area triangles appear on both the clipped surface and the cap.

The failure mode is worth stating precisely because its scope is
counter-intuitive. A single null-normal facet makes the Werner--Scheeres
evaluation return NaN at \emph{every} field point on the body, not merely
at points near the plane, because the facet's contribution enters the sum
for all field points. The symptom is therefore total rather than local,
and it cannot be mistaken for a numerical inaccuracy or worked around by
discarding nearby points.

The present implementation classifies vertices within $10^{-10}$ of the
bounding-box diagonal as lying on the plane, snaps them onto it exactly,
and never duplicates them; only strict sign changes generate an
intersection. On Itokawa that tolerance is $\SI{64}{\nano\metre}$. The row
reported at $\SI{0.1300}{\kilo\metre}$ in
Table~\ref{tab:itokawasweep} comes from the fixed implementation and is
not affected; the failure is recorded here because the plane at which it
occurs is one the sweep uses, so a reader reproducing this work without
the tolerance rule will meet it at the first plane. At the reference plane
of $\SI{0.1608}{\kilo\metre}$ no degenerate facet arises with or without
the rule.

The rule does not remove slivers, only null-normal facets. Slivers remain
and are harmless, which Section~\ref{sec:methods:tol} establishes by
measurement rather than by assertion.

\subsection{Regions and complements that are too small}
\label{sec:failures:small}

As the dividing plane advances the region shrinks, and the objective
becomes progressively less able to determine its density. The curvature
sensitivity interval on $\denshead$ widens from
$\pm\SI{130}{\kilo\gram\per\cubic\metre}$ at $f = 16.43\%$, which is
$4.6\%$ of the recovered value, to $\pm242$ at $f = 7.86\%$ ($6.2\%$) and
to $\pm355$ at $f = 5.00\%$, or $7.1\%$ of the recovered value. Extracted
from the scan grid rather than from a parabola, the intervals are
$-125/+125$ at the reference plane, $-200/+225$ at $f = 7.86\%$ and
$-350/+325$ at $f = 5.00\%$: asymmetric by one grid step at the two outer
planes and symmetric at the reference plane to the resolution of the scan.
% Resolved. The pair "-355/+1114" conflated two confidence levels rather
% than describing an asymmetry: at f = 5.00 % the grid gives -350/+325 at the
% 0.60 % level and -1125/+1100 at the 5.9 % level, and 1114 is the half-width
% at the latter. The level sets are indeed the roots of a quadratic and are
% asymmetric in general, but at these planes the asymmetry is one grid step.

A second observation at $\xsplit = \SI{0.220}{\kilo\metre}$ makes the
flatness concrete in an unexpected way. Head densities of $3700$ and
$\SI{4125}{\kilo\gram\per\cubic\metre}$, differing by $425$, return
dispersions differing by $4\times10^{-6}$. That is not by itself a
demonstration that the minimum is flat: two points placed nearly
symmetrically about a parabolic minimum return nearly equal values however
sharp the minimum is. Read as such a pair, however, the observation
locates the minimiser below the scan grid. With the local curvature fixed
by the $\pm242$ interval, equality to $4\times10^{-6}$ requires the two
displacements to differ by $\SI{1.4}{\kilo\gram\per\cubic\metre}$, placing
the stationary point at $\densR^{\ast} \simeq
\SI{3912}{\kilo\gram\per\cubic\metre}$. The grid reports $3900$, which is
$\SI{12}{\kilo\gram\per\cubic\metre}$ away and so within the half-step
bound claimed in Section~\ref{sec:methods:minimisation}---an inadvertent
confirmation of that bound, and a further argument for solving the
stationarity condition in closed form rather than scanning.

The opposite failure is the more dangerous of the two, and it is the one
this paper nearly reported. At $\xsplit = \SI{-10.76}{\kilo\metre}$ on
Eros the region is not small: it holds $88.44\%$ of the volume, and it is
the \emph{complement} that is small, at $11.56\%$. The mass constraint
then forces the complement density across an enormous range as $\densR$ is
scanned, from $\SI{19259}{\kilo\gram\per\cubic\metre}$ at
$\densR = 500$ to $\SI{127}{\kilo\gram\per\cubic\metre}$ at
$\densR = \SI{3000}{\kilo\gram\per\cubic\metre}$, so that most of the scan
describes a body that cannot exist. The recovered minimum has a curvature
interval of only $\pm\SI{15}{\kilo\gram\per\cubic\metre}$, three times
narrower than at the true saddle, and that apparent precision is the
symptom.

It is a coordinate artefact. The interval is narrow in $\densR$ because
the mass constraint makes $\densrest$ enormously sensitive to $\densR$:
by Equation~\eqref{eq:massbalance},
$\delta\densrest = [f/(1-f)]\,\delta\densR$, and at $f = 0.8844$ that
factor is $7.65$. Expressed in the quantity
Section~\ref{sec:constrains} concludes the method determines, the
interval is
$\delta\densrest = \SI{115}{\kilo\gram\per\cubic\metre}$, or $3.83\%$ of
the recovered $\SI{2996}{\kilo\gram\per\cubic\metre}$---the worst
determination anywhere in this study.

Table~\ref{tab:propagated} sets the four cases side by side, and the
ordering is the point of this subsection.

\begin{table}[htbp]
\centering
\small
\caption{The curvature sensitivity interval propagated through the mass
constraint. The interval on $\densR$ narrows as the region grows and
widens as it shrinks; the propagation factor $f/(1-f)$ moves the opposite
way and by more, so the interval on $\densrest$ is \emph{smallest} where
the region is smallest. A region that is too small therefore damages only
the quantity the method does not constrain, while a complement that is
too small damages the one it does.}
\label{tab:propagated}
\setlength{\tabcolsep}{10pt}
\begin{tabular}{llrrrr}
\toprule
Body & Plane (km) & $f$ (\%) & $\delta\densR$ & $f/(1-f)$
 & $\delta\densrest/\densrest$ \\
\midrule
Itokawa & $+0.1608$  & $16.43$ & $\pm130$        & $0.197$ & $1.38\%$ \\
Itokawa & $+0.2200$  & $7.86$  & $\pm242$        & $0.085$ & $1.11\%$ \\
Itokawa & $+0.2400$  & $5.00$  & $-355/+1114$    & $0.053$ & $-3.16\%/+1.01\%$ \\
Eros    & $-10.7619$ & $88.44$ & $\pm15$         & $7.654$ & $\pm3.83\%$ \\
\bottomrule
\end{tabular}
\end{table}

The practical rule that follows is stated in
Section~\ref{sec:guidelines}: where a body admits a choice, place the
\emph{smaller} of the two pieces in the region. It costs precision in a
quantity the method does not determine and buys it in the quantity it
does.

\subsection{Minima at a scan edge}
\label{sec:failures:edge}

A minimisation reports the smallest value it finds, and if the smallest
value is at the edge of the interval searched, it reports the edge. The
guard of Section~\ref{sec:methods:minimisation} exists because this
failure produces the most convincing-looking output in the catalogue.

The clear case is the Arrokoth neck model. A slab enclosing $1.0\%$ of the
volume drives $\densR$ to the scan floor of
$\SI{25}{\kilo\gram\per\cubic\metre}$ and reports an improvement of
$13.6\%$---four-fifths of the Itokawa figure, on a body that returns
$1.272\%$ under the head model. The guard flagged it and the result is
discarded. Widening the interval does not help here, because the limit is
physical rather than numerical: with $f = 0.01$ the mass constraint gives
$\densrest = \SI{403.8}{\kilo\gram\per\cubic\metre}$ and an excess mass of
only $-0.947\%$ of the total, so an extreme contrast is being purchased
with almost no mass.

That last figure suggests a diagnostic, which we put forward as a
candidate on the same footing as the mass-sensitivity test of
Section~\ref{sec:constrains:mass}. The ratio of the improvement to the
fractional excess mass---the dispersion reduction bought per unit of
claimed anomaly---is $14.4$ for the Arrokoth neck against $2.20$ for the
Itokawa head, $0.71$ for the Eros saddle, $0.66$ for the Itokawa neck,
$0.45$ for the misplaced Eros plane, $0.31$ for 67P and $0.26$ for the
Arrokoth head. The edge minimum exceeds the next highest value by a factor
of $6.5$ and is the only case that does. The diagnostic costs nothing, but
its limitation should be stated with it: it does not flag the misplaced
plane of Section~\ref{sec:failures:plane}, whose ratio is unremarkable.

The second case is the one this paper resolved rather than discarded. At
$\xsplit = \SI{0.220}{\kilo\metre}$ on Itokawa an early run used an upper
scan limit of $\SI{3900}{\kilo\gram\per\cubic\metre}$ and returned a
minimum at that limit. Widening the interval to
$\SI{8000}{\kilo\gram\per\cubic\metre}$ returned an interior minimum at
the same tabulated value. Two things should be said about that agreement.
The run with the lower ceiling predates the archived configuration, in
which the limit at that plane is $8000$, so it cannot be reproduced from
the released driver scripts and is described here only because the
diagnostic it illustrates is the point. And the agreement is one of grid
rather than of arithmetic: with no sub-grid refinement the minimiser is
reported to the nearest $\SI{25}{\kilo\gram\per\cubic\metre}$, and
Section~\ref{sec:failures:small} places the true stationary point near
$\SI{3912}{\kilo\gram\per\cubic\metre}$, which happens to round onto the
old boundary. A boundary minimum is a warning that the interval may be too
narrow; it is not a demonstration that it is.

A third consequence of the ceiling belongs here rather than in
Section~\ref{sec:plane}. Holding $\Dmass$ near its sweep mean, the ceiling
of $8000$ admits no region smaller than $f_{\min} = 2.56\%$
(Equation~\eqref{eq:capbinds}), which the sweep would reach at
$\xsplit \approx \SI{0.257}{\kilo\metre}$. Planes beyond that point are
not merely inadmissible on compositional grounds; they are outside what
the present scan configuration can resolve at all, and no statement about
the behaviour of $\potvarn$ there is supported.

\subsection{A misplaced plane}
\label{sec:failures:plane}

The plane is the one ingredient neither measured nor derived
(Section~\ref{sec:plane}), and misplacing it is the failure that produces
the most plausible output. On Eros the rejected selection rule of
Section~\ref{sec:methods:neck}---smallest interior half-width rather than
saddle---returns $\xsplit = \SI{-10.7619}{\kilo\metre}$, on a shoulder.

The output at that plane is worth setting out in full, because every
individual figure is unremarkable and the combination is wrong:
%
\begin{equation}
C = 0.876 \pm 0.039, \qquad
\improve = 5.526\%, \qquad
\Dmass/\Mtot = -12.32\%, \qquad
\dCOM = \SI{-214.8}{\metre} .
\label{eq:misplaced}
\end{equation}
%
The contrast excludes unity by $3.2$ of its own intervals. The excess mass
is larger in magnitude than the $7.7\%$ recovered on Itokawa. The
centre-of-mass offset is nearly four times that at the true saddle and of
the opposite sign. Nothing in the numbers marks them as spurious.

Four diagnostics were available and they do not agree, which is why this
subsection exists.

\emph{The envelope comparison catches it, narrowly.} At $5.526\%$ the
improvement lies inside the $0.9$--$5.9\%$ envelope of Paper~IV, but only
$6\%$ below its upper edge, and that edge is attained on a different body.
The comparison that matters is against the value the envelope takes on
Eros itself, which is $6.03\%$ over the ladder from $2{,}000$ to $42{,}000$
facets with quadric decimation. The misplaced plane reaches $0.92$ times that
value, so the envelope test rejects it, but with vertex clustering, which
Paper~IV used, the Eros envelope is $5.40\%$ and the same result would be
$1.02$ times it and would not be rejected: the margin is smaller than the
difference between two defensible ways of computing the envelope.

Two things should be said about the figure of $6.03\%$ itself. It exceeds
the upper edge of the $0.9$--$5.9\%$ range quoted throughout this paper, and
there is no contradiction in that: the published range was measured under
vertex clustering, whereas $6.03\%$ is measured under the quadric edge
collapse used here, and Appendix~\ref{app:envelope} gives both. Where a
body-specific value is available it supersedes the cross-body range, which is
attained on a different body in any case. The corresponding Itokawa value
under quadric collapse is $1.35\%$, so the $16.9\%$ improvement reported
there is $12.5$ times its own body's envelope, and the two comparisons are
made on the same footing.
% Resolved from paper4_recheck.json: the Eros-specific envelope is 6.03 %
% with quadric decimation and 5.40 % with vertex clustering. The misplaced
% plane reaches 5.526 %, so the diagnostic catches it under the first and
% not under the second; both values are now stated in the text.

\emph{Plane provenance catches it.} The plane was selected by a rule the
paper had already rejected on geometric grounds before any density was
computed, and the saddle prominences of Section~\ref{sec:methods:neck}
separate cleanly from $0.490$ on Arrokoth to $0.0004$ on Bennu. This is a
strong argument and it is not available to a reader handed only the
output.

\emph{The displacement statistic ranks it correctly, but weakly.} With
$\densrest = \SI{2996}{\kilo\gram\per\cubic\metre}$ and the propagated
interval of Table~\ref{tab:propagated},
%
\begin{equation}
\Sstat = \frac{\lvert\densbulk-\densrest\rvert}{\sigma(\densrest)}
       = \frac{328.6}{114.8} = 2.86 ,
\label{eq:smisplaced}
\end{equation}
%
against $5.8$ for the Itokawa reference solution. The artefact therefore
scores below the result, which is the right ordering; but it also scores
above all four genuine nulls, whose values run from $0.02$ on Bennu to
$0.68$ at the Eros saddle, so $\Sstat$ places it between the two classes
rather than with the nulls. It is a necessary condition and not a
sufficient one.

\emph{The amplification factor fails outright.} At this plane
$\ampfac = \densrest/(\densbulk-\densrest) = -9.12$, so
$\lvert\ampfac\rvert = 9.1$ against $12.0$ for Itokawa. By the ordering
established over five bodies in Section~\ref{sec:nulls:amplification}, in
which a smaller $\lvert\ampfac\rvert$ accompanies a stronger response,
this plane ranks as the strongest signal in the study. It is the one case
we have found that breaks that ordering, and it bounds what the
correlation reported there can be used for: $\ampfac$ orders bodies
evaluated at properly placed planes, and says nothing about whether the
plane is properly placed.

A fifth diagnostic was used during development and should be defined
before it is relied on. The fraction of the scan interval over which the
implied $\densrest$ remains within physically plausible bounds is
%
\begin{equation}
W = \frac{1}{\Delta\densR}\,\frac{1-f}{f}\,
    \left(\densrest^{\max}-\densrest^{\min}\right) ,
\label{eq:window}
\end{equation}
%
clipped to unity, where $\Delta\densR$ is the width of the scan. It is
computable from the mass constraint alone, before any gravity evaluation,
and it collapses as $f \to 1$, which is exactly the small-complement
regime of Section~\ref{sec:failures:small}. Evaluated with one rule for both bodies---each density confined to
$0$--$\SI{3540}{\kilo\gram\per\cubic\metre}$, the grain density of the
ordinary-chondrite analogue, and to the scan---it is $87.6\%$ for the
Itokawa reference plane, where the constraint binds only through
$\densR \le 3540$, and $10.3\%$ for the misplaced Eros plane, where the
small complement confines $\densR$ to $2554$--$\SI{3017}{\kilo\gram\per\cubic\metre}$
out of a scan of $500$--$5000$. The earlier figure of $90.0\%$ for Itokawa
was quoted against a slightly different upper bound and is superseded.
% Resolved: 10.3 % is reproduced exactly by applying the bounds to BOTH
% densities, which is what the working notes did; the ratio quoted in the
% query assumed the bounds act on rho_rest alone. Recomputed under the single
% rule stated in the text, Itokawa gives 87.6 % rather than 90.0 %.

\subsection{Facets that are not on the surface}
\label{sec:failures:buried}

The Arrokoth model of \citet{porter2024} is supplied as two unjoined
shells, geometrically in contact and topologically separate, with
$\chi = 4$ over two connected components. Four hundred and fifty-five
facets, carrying $0.943\%$ of the total area, have their centroids inside
the opposite shell, and the shells interpenetrate over $0.044\%$ of the
volume.

Those facets are not on the body's surface, and the relaxation argument of
\citet{richardson2014} does not apply to them: no regolith flows there and
no geopotential is levelled there. Including them in
Equation~\eqref{eq:dispersion} is a violation of the geometry, not an
approximation within it.

The consequence is larger than the area warrants. At an assumed bulk
density of $\SI{250}{\kilo\gram\per\cubic\metre}$, retaining the buried
facets returns a contrast of $1.062$ and excluding them returns $0.9999$:
under one per cent of the area moves the answer by $6.2\%$ and creates an
apparent density excess where there is none. Across the assumed-density
range the same decision moves the crossover---the assumed bulk density at
which the contrast passes through unity---from between $235$ and
$\SI{250}{\kilo\gram\per\cubic\metre}$, and within one kilogram per cubic
metre of $250$ by interpolation between those two runs, to approximately
$\SI{345}{\kilo\gram\per\cubic\metre}$
(Section~\ref{sec:nulls:arrokoth}).


All Arrokoth results reported in this paper exclude them, with $\Ubar$ and
$\sum_i A_i$ recomputed over the retained set rather than terms merely
dropped from the numerator. The check that finds them is a containment
test of each facet centroid against the other components, and it is
required only because the source model has more than one component---which
the per-component Euler criterion of Section~\ref{sec:methods:verify}
reports, and a single global $\chi = 2$ test would not.

\subsection{What the catalogue implies}
\label{sec:failures:summary}

Three of the seven failures announce themselves: the degenerate facets
produce NaN, the edge minima are caught by the boundary guard, and the
region-too-small case declares itself through a curvature interval that
reaches $22\%$ of the recovered value. Four do not.

Of the four silent ones, the wrong solid and the buried facets are
geometric and are caught by checks on the region's \emph{definition}
rather than on its self-consistency---the vertex predicate, the
per-component Euler criterion, the containment test, and the mass-balance
audit. The small complement is caught only by propagating the curvature
interval through the mass constraint before reading it, since in the
scanned coordinate it presents as unusual precision. The misplaced plane
is caught only by the provenance of the plane and, marginally, by
comparison against the body's own resolution envelope; two of the four
statistical diagnostics available rank it below the genuine result, and
one ranks it above.

The general shape of the catalogue is that no single number distinguishes
a result from an artefact. This is why Section~\ref{sec:nulls:misplaced}
states a three-part criterion rather than a threshold, why
Section~\ref{sec:guidelines} asks for the plane, its provenance and the
propagated interval to be reported alongside any density, and why the
quantity we recommend reporting is $\densrest$ rather than $\densR$: of
the seven failures, one degrades $\densrest$ and six leave it intact or
destroy the calculation outright, whereas five of the seven corrupt
$\densR$ while leaving the output superficially reasonable.
